\scanpage{9}
% archive sheet number: 8
\[
\mathrm{R}f_{\handwritten{!}}(\handwritten{f^*}G(d-d'))\handwritten{\longrightarrow}G\handwritten{[}2(d'-d)\handwritten{]}.
\tag*{(4.3 bis)}
\]
Arrivé à ce point, je constate que c'est par cette flèche\inserted{(4.3)} (plus raisonnable que (4.2)) qu'il aurait fallu commencer, en ajoutant encore que si $p$ est également lisse et est de \unsure{dim,|dim;} rel. $d$ partout, \deleted{alors (\unsure{4.3|4.2})}\inserted{et si $G$ est localement (sur $Y_{\text{ét}}$) isomorphe à l'image inverse d'un complexe sur $S$, alors (4.3)} s'identifie à l'homomorphisme d'adjonction qu'on pense. Il y aurait lieu alors de donner les exemples (4.2.1) à (4.2.8) \deleted{\unsure{du|de}} en partant de (4.3) plutôt que de (4.2). On peut laisser tomber (4.2.1), et se borner à dire que les homomorphismes déduits des homomorphismes envisagés par application des foncteurs $\mathrm{H}^i$ sont également affublés du nom ``Gysin'' ou ``trace''.

\unsure{Une|Upe} perplexité: dans le contexte actuel, la notation $f_{\handwritten{*}}$, qui dénote à la fois un foncteur (image directe) et un homomorphisme de foncteurs (Gysin) ne semble pas bien fameux. Je propose de donner la préférence ici à la notation $\mathrm{Tr}_f$, en le disant au moment ou on signale également la notation $f_{\handwritten{*}}$.

Etant donné le nombre d'exemples parmi les variantes à Gysin, il est clair qu'il ne sera pas possible de donner un formulaire qui donne\inserted{rait} \deleted{immédiatement} sans déductions du tout, tous les cas qu'on rencontrera. Il faudra se borner à donner un formulaire directement en termes des homomorphismes de Gysin fondamentaux sous la forme (4.3), et se borner à donner comme exemple la forme \unsure{qae|que} prend ce formulaire\handwritten{``}après application d'un foncteur $\mathrm{R}g_*$\handwritten{''} dans le cas $f$ propre, disons.

Voici des points qui doivent certainement figurer au formulaire:

\margin{\handwritten{descend}}
\deleted{\unsure{a}}\handwritten{d}) Compatibilité avec la formule de projection: Si on a un complexe de modules $G'$ sur $Y$ et si on tensorise (4.3) par $G'$, en utilisant la formule de projection de XVII ... , on trouve Gysin pour $G\overset{\mathrm{L}}{\otimes}G'$. Ceci \deleted{montre}\inserted{\handwritten{\unsure{contient}}} en particulier que, via la formule de projection,
